
\chapter{An Introduction to Rietveld Fitting}

In this book, I introduce a number of techniques for analysis of powder diffraction data, concentrating on most, but not all, of the capabilities within the GSAS-II software package. 
I devote the most attention to Rietveld analysis\index{Rietveld analysis|textbf}, which is a process where a model is fitted directly to a powder diffraction pattern\index{powder diffraction}. This method was initially developed to determine the arrangement of atoms in materials, or as this process is commonly called, determining the crystal structure. 

Crystal structure determination was done with single-crystals as far back as 1913, and early powder diffraction work was also used for this later in that same decade. Structure determination from single crystals prospered greatly in the subsequent half-century, as instruments improved, computational model-fitting was adopted, and structure solution techniques were developed. Computerization helped considerably, but it could be argued that single-crystal work has advanced even more in the last few decades. 
Powder diffraction crystallography, on the other hand, was largely left behind in that first half-century, at least until Hugo Rietveld\index{Rietveld, H.} and some of his coworkers at the Netherlands' Petten\index{Petten, Netherlands} High Flux Nuclear Reactor pursued the problem in the 1960's. 

At that time neutron single-crystal measurements required crystals the size of a finger-tip (which were quite rare) and weeks, if not months, of continuous data collection. So, there was a plethora of samples where only powder diffraction could be measured. It was true then, and is still true now, that when single-crystal data are available for a material, the single-crystal measurement will provide much more structural detail. However, the single-crystal data needs to be from a material with the desired composition and where the measurements are under the conditions of interest. There are many materials where single-crystals cannot be formed with the desired chemical or physical properties. Or, measurements are not possible under the desired conditions: in some cases one wants to learn about processes that occur under conditions that degrade single-crystals to a polycrystalline material. Thus, while I consider powder diffraction as a second-choice technique, it is often the only choice. With powder diffraction, it becomes relatively easy to make a combined fit using multiple types of data, such as x-ray and neutron, or with data collected near x-ray absorption edges (known as resonant scattering or anomalous dispersion measurements). In fact, x-ray single-crystal measurements results can often be enhanced, when combined with powder diffraction measurements with neutrons or other x-ray wavelengths.

Further, in the subsequent decades since Hugo's work, we have found many other applications for this type of fitting. It provides information not only about atomic arrangements (sub-nanoscale structure), but also about on larger-scales than atomistic information (microstructure) such as crystallite size and texture. When materials are mixtures, and they commonly are, either by intent, as mixtures can improve properties, or because it not known how to create the pure phase, powder diffraction can provide quantification of the components of the mixture. Powder diffraction measurements can be quick. In fact, as far as I am aware the world's record has been a complete measurement made 100 picoseconds.  Also, measurements can be made under \textit{in situ} or operando conditions. Due to the speed and simplicity of the measurement, it is now common to perform parametric experiments where even thousands of complete datasets are collected on a sample as conditions are changed. It is no surprise that when a neutron or synchrotron center builds a powder diffractometer, while not the most sexy cutting-edge of techniques at the institution, these instruments go on to become some of the most in-demand and most productive instruments at the facility. 

Sometimes the term full-pattern fitting is used to describe this type of work, rather than Rietveld analysis. I prefer to honor Hugo Rietveld\index{Rietveld, H.}. He was not the only person to work on ideas for fitting powder diffraction patterns, but was the first to put the technique into practice. In developing such a code, he had to overcome a number of technical difficulties with, by modern standards, tiny and slow computers. After showing that the technique was possible and of value, he rewrote his code in a more accessible computer language and freely distributed his code to anyone who wanted a copy. A few decades after Hugo’s work was completed, there were a number of Rietveld codes in use, but almost all contained some code that Hugo had developed. I have written a short history of Rietveld analysis as chapter 4.7 in the International Tables for Crystallography, Volume H, for those who might want to learn more and this provides citations to the original literature. 

So, what is the purpose of this book? Herein, I want to present the important concepts needed for working with powder diffraction analysis, but with a modern and hand-on approach. This is work that requires computer software, so my goal is to present how and why these analyses are done, accompanied with theory when appropriate, but then present how to do that task using the widely-used and freely available crystallographic data analysis package GSAS-II\index{software!GSAS-II}. I think this is the most effective way to learn to become proficient in these analysis. 

I am one of the two principal authors of GSAS-II and I am quite proud of the work done by my partner in its creation, Robert B. Von Dreele\index{Von Dreele, R. B.} and I. We continue to work to make GSAS-II the best that we can, but I do not want to make the assertion that GSAS-II is always the best tool for the job. There are a number of fine packages that include Rietveld analysis, including Topas\index{software!Topas}, FullProf\index{software!FullProf}, JANA\index{software!JANA}, Rietan\index{software!Rietan} and MAUD\index{software!MAUD}. Each offers some capability that GSAS-II does not. I am not sufficiently knowledgeable to introduce use of any of those packages here, so I will stick to GSAS-II, but not because I think it is the only program that can do the things I present. Nonetheless, I do think that GSAS-II has the widest range of features; it is the most modern, started decades after many of the others; we work at making GSAS-II easily accessible to people getting started in the field; it runs on all types of computers in common use; and, you can’t beat the price. 

However, Bob and I do not recommend the use of our previous software collaboration, GSAS/EXPGUI\index{software!GSAS}\index{software!EXPGUI}. This software has been superseded in every way by GSAS-II. There will be no future updates to GSAS or EXPGUI. As of today, GSAS/EXPGUI does run on Windows, but only with great difficulty on other modern platforms. GSAS was a marvel in its day. It had a very modern (non-graphical) user interface, at least in the 1980's. GSAS was probably the first Rietveld program to have been written that did not contain any of Hugo’s original code. EXPGUI added a then-modern (1990's) graphical user interface (GUI) to some of GSAS. Nonetheless, we can't keep both packages alive, so it is time to say goodbye to GSAS/EXPGUI. Even the compilers/scripting packages needed for them are obsolete.

So what is Rietveld analysis about? As mentioned, it is the fitting of an as-measured powder diffraction pattern. In single-crystal fitting, one only needs to fit to the observed structure factor intensities (which, by the way, are not directly measured, but must be extracted as part of the single-crystal experimental process). In both, one optimizes details of the atomistic structure of the phase to best-fit the observed intensities. However, with powder diffraction, it now may be necessary to fit more than one structure, as a material may contain many phases. Additionally, in Rietveld fitting, there are three more major aspects in the as-measured powder diffraction pattern that must be fit: 

\begin{itemize}

    \item There will always be non-Bragg scattering intensity present. We call this background. In the single-crystal measurement, this is automatically removed (one hopes -- well), but in powder diffraction the background levels are more substantial, can vary very significantly across the pattern, and must be modeled point by point. 

    \item  Peak positions are determined by the lattice parameters, so these need to be fit as well. Fortunately, powder diffraction is both extremely accurate and precise in this measurement – frequently an order of magnitude better than single crystal measurements. 
    
    \item In order to fit the as-measured powder diffraction pattern point-by-point, we must account not only for the intensity of the reflections from the structure, but also the way that these reflections are broadened (in single-crystal work, one just integrates over the breadth.) This broadening arises from both the sample and the instrument and is modeled by the peak shape function. When multiple data types are combined, the breadths need to be treated differently for each phase, as well as for each measurement when more than one set of experimental data is used. 

\end{itemize}

In later sections of this book I will discuss how each aspect of the four aspects modeling is done, but it should be clear that any metrics for how well the data are fitted will reflect the overall quality with respect to all four aspects of the fit. In Chapter \ref{Chap:rfactors}, I will discuss these metrics in detail, but for now let us concentrate on the quantity called the R-factor. 
Powder diffraction R-factors are poor at reflecting the quality of a structural fit since R-factor indicates how well all four aspects of the fit. If a powder diffraction pattern contains a large amount of background, and that is well-fit, but the structure is not, the R-factor may be low. On the other hand, if minor aspects of the background are poorly fit, a high R-factor may result, even though the structural fit is excellent. In contrast, a single-crystal R-factor only measures how well the crystal structure reproduces the observed intensities, so it is a more clear indication of quality – though often overemphasized. 

As a professional, employed at neutron and synchrotron sources for most of my career, as well as through Rietveld software support, I have interacted with many people who are attempting to get started with Rietveld analysis. These are typically very bright and well-educated young scientists and engineers; nearly all are a pleasure to work with. However, it is more frequent than I would prefer that I encounter someone who does not understand key aspects of the what the measurements reflect and what the software is doing. It makes it hard for them to know how to do things correctly and if the results they have obtained are meaningful. I strongly believe that working with Rietveld refinement requires a good understanding of diffraction, symmetry and powder diffraction. I think modern scientists who want to perform analysis of powder diffraction are at a severe disadvantage without mastery of these concepts.

I may be prejudiced, since my introduction to crystallography was through single crystal diffraction – and sufficiently long ago so that the methods were then certainly far from “black box” techniques – but I believe having this background allowed me to come to powder diffraction with a good understanding of many basic crystallographic concepts. While I do not want to write yet another textbook on basic crystallographic concepts – there already are plenty, I think a valuable place to start is to review foundational concepts and I will try to show how they are interrelated. For this reason I will start this book with three chapters labeled as “prerequisites.” My hope is that this material will largely be a review to you, though I do also hope that on occasion I may present some of this in a new way. Should any of this material be unclear to you as you read these prerequisite chapters, I suggest you make an immediate detour and read about the topic(s) in a basic crystallography text. The subsequent chapters will then be much easier to understand and you will be much better prepared to perform Rietveld analysis and judge the quality of your results. 